\section{Análisis de incertidumbre frecuentista}
\label{sec:incertidumbre}

La proporción observada de malos créditos describe la composición de la
muestra, pero no expresa por sí sola la precisión de la estimación. Para
cuantificar esta incertidumbre se construyeron dos intervalos de confianza
del \(95\,\%\) para la tasa de mal crédito: un intervalo analítico de
Wilson y un intervalo bootstrap percentil.

Como análisis complementario, se estimó también un intervalo basado en
la distribución \(t\) de Student para la duración media del crédito. Esta
segunda estimación permite distinguir la incertidumbre asociada con una
proporción binomial de la incertidumbre correspondiente a la media de una
variable cuantitativa.

Los resultados completos se generan en el notebook y se exportan en:

\begin{center}
    \texttt{resultados/tablas/tabla\_intervalos\_confianza.csv}
\end{center}

\subsection{Proporción observada de malos créditos}

Sea \(Y_i\) la variable binaria que identifica un mal crédito:

\[
Y_i =
\begin{cases}
1, & \text{si la observación corresponde a mal riesgo},\\
0, & \text{si corresponde a buen riesgo}.
\end{cases}
\]

Bajo un modelo binomial:

\[
X=\sum_{i=1}^{n}Y_i
\sim
\operatorname{Binomial}(n,p),
\]

donde \(p\) representa la proporción poblacional desconocida de malos
créditos.

En la muestra analizada:

\[
n=1000,
\qquad
x=300,
\]

por lo que la estimación puntual es:

\[
\widehat{p}
=
\frac{x}{n}
=
\frac{300}{1000}
=
0.300.
\]

La tasa observada de mal crédito es, por tanto, del \(30.0\,\%\).

\subsection{Intervalo analítico de Wilson}

Para estimar \(p\) se utilizó el intervalo de Wilson. A diferencia del
intervalo de Wald, este procedimiento incorpora una corrección tanto en
el centro como en el margen del intervalo y suele presentar una cobertura
más estable para proporciones binomiales.

Para un nivel de confianza \(1-\alpha\), el intervalo se define como:

\[
\operatorname{IC}_{\mathrm{Wilson}}
=
\frac{
    \widehat{p}
    +
    \dfrac{z_{1-\alpha/2}^{2}}{2n}
    \pm
    z_{1-\alpha/2}
    \sqrt{
        \dfrac{\widehat{p}(1-\widehat{p})}{n}
        +
        \dfrac{z_{1-\alpha/2}^{2}}{4n^{2}}
    }
}{
    1+\dfrac{z_{1-\alpha/2}^{2}}{n}
}.
\]

Para un nivel de confianza del \(95\,\%\):

\[
\alpha=0.05,
\qquad
z_{1-\alpha/2}\approx 1.96.
\]

El intervalo obtenido fue:

\[
\operatorname{IC}_{95\,\%,\,\mathrm{Wilson}}
=
[0.2724,\;0.3291].
\]

En términos porcentuales:

\[
27.24\,\%
\leq
p
\leq
32.91\,\%.
\]

La amplitud total del intervalo es de aproximadamente \(5.67\) puntos
porcentuales. De forma equivalente, el margen alrededor de la estimación
central es cercano a \(2.84\) puntos porcentuales.

El cálculo se implementó explícitamente en el notebook y se verificó
mediante la función \texttt{proportion\_confint()} de
\texttt{statsmodels}. Ambas implementaciones produjeron los mismos
límites dentro de la tolerancia numérica utilizada.

\subsection{Intervalo bootstrap percentil}

Como método de remuestreo se utilizó un bootstrap no paramétrico. El
procedimiento genera \(B=10\,000\) muestras con reemplazo, cada una de
tamaño \(n=1000\), a partir de los registros observados.

Para cada réplica \(b\) se calculó:

\[
\widehat{p}^{*(b)}
=
\frac{1}{n}
\sum_{i=1}^{n}
Y_i^{*(b)}.
\]

La colección:

\[
\widehat{p}^{*(1)},
\widehat{p}^{*(2)},
\ldots,
\widehat{p}^{*(B)}
\]

aproxima la distribución muestral de la proporción. El intervalo
percentil se construye mediante los cuantiles \(0.025\) y \(0.975\):

\[
\operatorname{IC}_{\mathrm{bootstrap}}
=
\left[
    Q_{0.025}\left(\widehat{p}^{*}\right),
    Q_{0.975}\left(\widehat{p}^{*}\right)
\right].
\]

Con la semilla reproducible \(\texttt{SEED}=2026\), el resultado fue:

\[
\operatorname{IC}_{95\,\%,\,\mathrm{bootstrap}}
=
[0.2720,\;0.3290].
\]

En términos porcentuales:

\[
27.20\,\%
\leq
p
\leq
32.90\,\%.
\]

El error estándar estimado a partir de las réplicas bootstrap fue:

\[
\widehat{\operatorname{SE}}_{\mathrm{bootstrap}}
=
0.0145.
\]

La amplitud total del intervalo bootstrap es de \(5.70\) puntos
porcentuales, prácticamente igual a la obtenida mediante Wilson.

\subsection{Comparación de los intervalos para la proporción}

La Tabla~\ref{tab:intervalos-proporcion} resume los resultados.

\begin{table}[H]
    \centering
    \caption{Intervalos de confianza del \(95\,\%\) para la proporción
    de malos créditos.}
    \label{tab:intervalos-proporcion}

    \begin{tabular}{
        @{}
        l
        c
        c
        c
        c
        @{}
    }
        \toprule
        Método
        & Estimación
        & Límite inferior
        & Límite superior
        & Amplitud \\
        \midrule

        Wilson
        & \num{0.3000}
        & \num{0.2724}
        & \num{0.3291}
        & \num{0.0567} \\

        Bootstrap percentil
        & \num{0.3000}
        & \num{0.2720}
        & \num{0.3290}
        & \num{0.0570} \\

        \bottomrule
    \end{tabular}
\end{table}

Los dos procedimientos producen resultados casi idénticos. La diferencia
entre sus límites inferiores es de solo \(0.0004\), mientras que la
diferencia entre los límites superiores es de \(0.0001\).

Esta concordancia se explica por el tamaño de la muestra y porque la
proporción observada no se encuentra cerca de los extremos cero o uno.
En estas condiciones, la distribución muestral de
\(\widehat{p}\) es aproximadamente simétrica y el remuestreo reproduce
de manera cercana la incertidumbre capturada por el intervalo analítico.

\begin{hallazgo}
La tasa de mal crédito se estima en \(30.0\,\%\), con límites cercanos a
\(27.2\,\%\) y \(32.9\,\%\). La concordancia entre Wilson y bootstrap
indica que la estimación es estable frente a los dos procedimientos de
construcción considerados.
\end{hallazgo}

\subsection{Intervalo para la duración media}

También se construyó un intervalo de confianza para la media poblacional
de la duración del crédito. Sea:

\[
\overline{x}
=
\frac{1}{n}
\sum_{i=1}^{n}x_i
\]

la duración media muestral y \(s\) la desviación estándar muestral. El
intervalo basado en la distribución \(t\) de Student es:

\[
\overline{x}
\pm
t_{1-\alpha/2,n-1}
\frac{s}{\sqrt{n}}.
\]

La estimación observada fue:

\[
\overline{x}
=
20.903\text{ meses}.
\]

El intervalo del \(95\,\%\) fue:

\[
\operatorname{IC}_{95\,\%,\,t}
=
[20.1547,\;21.6513]
\text{ meses}.
\]

La amplitud total es de aproximadamente \(1.50\) meses. Este resultado
indica que la media de la duración se estima con una precisión
relativamente alta dentro de la población representada por la muestra.

Debe distinguirse esta incertidumbre de la variabilidad entre créditos.
El intervalo estrecho corresponde a la precisión de la media estimada;
no significa que las duraciones individuales se encuentren próximas a
\(20.9\) meses. Como se observó en la estadística descriptiva, los
contratos varían entre \(4\) y \(72\) meses.

\begin{table}[H]
    \centering
    \caption{Resumen de los intervalos de confianza construidos.}
    \label{tab:resumen-intervalos}

    \begin{tabularx}{0.94\textwidth}{
        @{}
        >{\raggedright\arraybackslash}p{3.2cm}
        >{\raggedright\arraybackslash}p{3.4cm}
        >{\centering\arraybackslash}p{2.0cm}
        >{\centering\arraybackslash}p{2.0cm}
        X
        @{}
    }
        \toprule
        Parámetro
        & Método
        & Estimación
        & IC del \(95\,\%\)
        & Unidad \\
        \midrule

        Proporción de mal crédito
        & Wilson
        & \num{0.300}
        & \([0.2724,\;0.3291]\)
        & Proporción \\

        Proporción de mal crédito
        & Bootstrap percentil,
          \(B=\num{10000}\)
        & \num{0.300}
        & \([0.2720,\;0.3290]\)
        & Proporción \\

        Media de la duración
        & \(t\) de Student
        & \num{20.903}
        & \([20.1547,\;21.6513]\)
        & Meses \\

        \bottomrule
    \end{tabularx}
\end{table}

\subsection{Interpretación frecuentista}

Los intervalos no deben interpretarse como una probabilidad posterior
sobre un parámetro fijo. Una vez calculado el intervalo, el parámetro
poblacional se encuentra dentro o fuera de sus límites; no se asigna una
probabilidad del \(95\,\%\) a ese evento dentro del enfoque
frecuentista.

La interpretación se refiere al procedimiento: si se extrajeran muchas
muestras comparables y se construyera un intervalo mediante el mismo
método en cada una, aproximadamente el \(95\,\%\) de esos intervalos
contendría el verdadero parámetro poblacional.

Esta interpretación se diferencia del intervalo creíble bayesiano que se
presentará en la Sección~\ref{sec:bayes}. En el enfoque bayesiano, la
incertidumbre se expresa directamente mediante una distribución
posterior para \(p\).

\subsection{Supuestos y alcance}

El intervalo de Wilson se apoya en un modelo binomial y supone que las
observaciones representan ensayos comparables e independientes. El
bootstrap utilizado también remuestrea registros individuales, por lo que
mantiene implícitamente una estructura de independencia entre las filas.

Además, ambos métodos cuantifican la incertidumbre muestral condicionada
a la composición del dataset. No corrigen posibles sesgos derivados del
proceso de selección original, la antigüedad de los datos ni la falta de
representatividad respecto de una cartera actual.

Por ello, el intervalo:

\[
[27.2\,\%,\;32.9\,\%]
\]

describe la precisión de la tasa dentro de una población comparable con
la muestra analizada. No debe interpretarse como una estimación directa
de la tasa de incumplimiento de una entidad financiera peruana ni de una
cartera contemporánea sin una validación externa adicional.

La cercanía entre Wilson y bootstrap aporta evidencia de estabilidad
numérica, pero no elimina las limitaciones de validez externa del
dataset.

% Contenido previsto:
% - IC de Wilson al 95 % para la tasa de mal crédito.
% - IC bootstrap.
% - Procedimiento de remuestreo, semilla y número de réplicas.
% - Interpretación de precisión e incertidumbre.
